\emph{Översikt}
Den här modulen är föreläsningens andra halva.  I den första halvan
(\emph{Algoritmiskt tänkande}) förfinade vi en vardagsalgoritm ända ner
till ett pythonprogram --- men vi körde det aldrig.  Här tar vi steget:
vad \emph{är} ett program, vad \emph{är} ett programmeringsspråk, och vad
händer i datorn när ni skriver \mintinline{bash}{python3 hello.py} i
terminalen?  Vi går från processorn, som bara kan utföra
maskinkod, till samma lilla hälsning skriven i tre språk --- Python, C++
och skalet --- och urskiljer ur kontrasten vad ett programmeringsspråk är
och varför något alltid måste stå emellan er text och processorn.  Sedan
kör vi Python på två sätt, interaktivt och som en sparad fil, för att
skilja \emph{filen} från det \emph{körande programmet} och redigeraren
från språket.  Modulen slutar där er laboration börjar: vid det första
felmeddelandet.

\emph{Lärandemål}
Efter denna modul ska studenten kunna:
% Ingen av kursens Lo-koder träffar den här modulen direkt: den är
% verktygs- och begreppssteget som resten av Pythondelen står på.
% Kopplingarna nedan anges därför som förkunskap där de är sådana.
\begin{restatable}{lo}{HelloLOProgram}\label{HelloLOProgram}%
  Förklara vad ett program och ett programmeringsspråk är: att processorn
  bara utför maskinkod, och att det vi skriver därför alltid måste
  översättas till maskinkod eller läsas av ett program som redan är det.
\end{restatable}
\begin{restatable}{lo}{HelloLOExecution}\label{HelloLOExecution}%
  Förklara vad som händer när Python kör en fil: tolken utför en sats i
  taget, uppifrån och ner, och gör bara det som står --- den gissar
  ingenting. % förkunskap för Lo04 (styrstrukturer)
\end{restatable}
\begin{restatable}{lo}{HelloLOTools}\label{HelloLOTools}%
  Skilja på textredigeraren, filen på disken och det körande programmet,
  och köra samma pythonprogram på båda sätten: interaktivt och som en
  sparad fil startad från terminalen.
\end{restatable}
\begin{restatable}{lo}{HelloLOWrite}\label{HelloLOWrite}%
  Skriva, spara och köra ett litet pythonprogram som skriver ut text,
  och byta ut det utskrivna mot en variabel och mot en funktion.
  % Lo05 (utmatning), Lo07 (identifierarnamn)
\end{restatable}
\begin{restatable}{lo}{HelloLOError}\label{HelloLOError}%
  Läsa Pythons felmeddelande: hitta filen och raden det pekar på, avgöra
  vad som hann köras innan felet, och rätta felet.
  % förkunskap för Lo11 (granska andras program)
\end{restatable}

\emph{Förkunskaper}
Modulen förutsätter inga tidigare programmeringskunskaper.  Studenten bör
ha sett föreläsningens första halva, \emph{Algoritmiskt tänkande}:
programmet \texttt{pannkakor.py} därifrån är det program vi kör här, och
notationen med namngivna steg införs där.

\emph{Läsning}
Terminalen och filsystemet lär ni er av materialet från
terminalkursen DD1301 överst i veckans modul (från \emph{Briefly on
interfaces} till \emph{Choosing an editor}).  Modulen är underlaget för
\emph{Laboration 0: kom igång med Hello World}.
